\section{Assessment Methodology}
\label{sec:methods}
The assessment examines the structure of the extracted events, their consistency with generated descriptions, and their use in the reported descriptive analyses. It combines structural checks, the existing author audit, and calculations with explicitly defined measures.

The inputs are frozen texts, reports, CSVs, and archived author judgments. The descriptive workflow runs offline with Python's standard library. Activity discovery uses PM4Py~\cite{berti2023pm4py} with the archived miner and noise-threshold settings. All calculations use the released records as their input.

\begin{table*}[!t]\begin{adjustwidth}{-\extralength}{0cm}
\caption{Evaluation questions, required references, and evidence in this study. Successful computation and correct operational interpretation are separate claims.}
\label{tab:evaluation_scope}\footnotesize
\begin{tabularx}{\linewidth}{@{}LLL@{}}\toprule
Question & Reference needed & Present evidence and scope \\\midrule
Are visual actions correct? & Independent annotations of the video & Not measured; generated transcripts cannot detect their own visual errors \\
Are labels consistent with text? & A reference interpretation of the text & Existing author judgments for 42 events; label membership checked for all rows \\
Are boundaries accurate? & Independently defined action boundaries & Agreement with generated transcript endpoints only; the five-second tolerance is a matching rule \\
Can the log support process-mining computations? & Defined case rules, measures, and completed calculations & Discovery and descriptive outputs demonstrate computation on the inferred representation \\
Are operational conclusions correct? & Valid measurements and evidence relevant to the conclusion & Traceable category descriptions and hypotheses; no independent production outcomes or intervention evidence \\
\bottomrule\end{tabularx}
\end{adjustwidth}\end{table*}

\subsection{Evaluation questions and reference data}
Table~\ref{tab:evaluation_scope} distinguishes the objects of evaluation. Our generated transcript checks assess an intermediate transformation, while our downstream calculations demonstrate uses of its output. These assessments support different conclusions, so their accuracy, fitness, and precision values should not be ranked as if they measured the same property.

\subsection{Metrics, denominators, and aggregation}
\label{sec:metrics}
Table~\ref{tab:metrics} defines each measure, its denominator, and how recordings or workers contribute to the result.

\emph{Coverage} is the fraction of the parsed timestamp span occupied by extracted events. For video $v$, let $E_v$ be its events, $B_v$ the transcript endpoints recognized by the original parser, $T_v^{\max}$ their maximum, and $M_v$ the latest event endpoint. Factory coverage is:
\begin{equation}
 C_v=\left|\bigcup_{e\in E_v}[t_s,t_e)\right|,\qquad S_v=\max(T_v^{\max},M_v),\qquad
 \mathrm{coverage}_f=\frac{\sum_{v\in f}C_v}{\sum_{v\in f}S_v}.
\end{equation}
$C_v$ is occupied time in seconds, with overlaps counted once; the vertical bars denote interval duration. $S_v$ runs from zero to the latest transcript or event endpoint. Dividing factory totals weights recordings by this span. If no transcript endpoint is recognized, $M_v$ supplies the denominator. These spans come from recorded timestamps; actual video durations are unavailable. We retain the original parser's results and release the broader parser's results separately.

% Declare the metric table early for top placement on page 17.
\FloatBarrier
\begin{table*}[!t]\begin{adjustwidth}{-\extralength}{0cm}
\caption{Metric definitions. $E_f^*$ is the retained set under $M_v\leq3600$; $E_w^*$ is its worker subset. Summed event-hours are $H_w=\sum_{e\in E_w^*}d_e/3600$. $R_w^{VA}$ is the worker's list of VA run durations, converted to minutes. Boundaries in recordings with no recognized generated transcript endpoints count as unmatched.}
\label{tab:metrics}\footnotesize
\begin{tabularx}{\linewidth}{@{}p{0.21\linewidth}p{0.39\linewidth}L@{}}\toprule
Quantity & Definition & Unit, aggregation, exclusions \\\midrule
Vocabulary membership & $\#\{e:a\in A_f,\ p\in P_f\}/\#E_f$ & Fraction of all event rows \\
Vocabulary use & Distinct assigned labels divided by permitted labels & Separate activity and process lists, per factory \\
Coverage diagnostic & $\sum C_v/\sum S_v$ & Fraction of parsed endpoint span; all logs \\
Boundary agreement & $\sum_{e,b\in\{t_s,t_e\}}\mathbf{1}[\min_{t\in B_v}|b-t|\leq5]/(2\#E)$ & Fraction of boundaries; all logs; original parser \\
Category share & $q_{fc}=\sum_{e\in E_f^*}d_e\mathbf{1}[g_f(p_e)=c]/\sum_{e\in E_f^*}d_e$ & Fraction of summed event duration; overlap counted in each row \\
Worker VA share & Same ratio restricted to $E_w^*$ and category VA & Eligible workers require $H_w\geq0.5$; factory median weights workers equally \\
VA transition rate & $I_w/H_w$, where $I_w$ counts VA-to-non-VA transitions in concatenated run lists & Transitions per summed event-hour; factory median \\
Mean VA run & $\overline R_w=\sum R_w^{VA}/\#R_w^{VA}$ & Minutes; table reports median across eligible workers of these means \\
MH/TR count and rate & Number of retained MH or TR rows; count divided by summed event-hours & Rows, not merged episodes; duration-weighted factory denominator \\
Dominant process & $\arg\max_p\sum_{e\in E_w^*:p_e=p}d_e$ & Groups workers; comparisons then use all their retained events \\
Self-loop share & Equal adjacent activity labels divided by all adjacent pairs & Within video cases only; all logs \\
\bottomrule\end{tabularx}
\end{adjustwidth}\end{table*}

\emph{Boundary agreement} is the fraction of event starts and ends within five seconds of a recognized transcript endpoint. Each event contributes two boundaries. If a recording has no recognized endpoints, its boundaries count as unmatched. Several boundaries may match one endpoint. This measures timestamp agreement without establishing action accuracy or validating an internal split.

The process map assigns events to \emph{value-adding production work} (VA), \emph{material handling and staging} (MH), \emph{transport} (TR), \emph{quality inspection} (QI), \emph{cleaning and maintenance} (CM), \emph{rework} (RW), and \emph{documentation} (DOC). VA includes packaging and completion work under the authors' convention; it does not directly measure customer value. Defect patching is RW because its descriptions refer to surface repairs, without establishing avoidability or failure against a quality standard. Work outside VA may still be necessary.

Of the 113 factory-specific process entries (108 distinct strings), 112 are used and mapped. The unused factory \textbf{005} label \emph{motor assembly finishing} has no category assignment. Appendix~\ref{supp-app:mapping} of the Additional Material lists the full map.

\emph{Category shares} divide a category's summed event duration by that of all retained events. Retention requires $M_v\leq3600$; $E_f^*$ and $E_w^*$ denote the retained factory and worker events. The process label determines the category, and each row contributes its full duration, including overlaps.

\emph{Worker summaries} use the same retained events. Worker event-hours $H_w$ are summed durations divided by 3600. Factory medians give eligible workers equal weight. The \emph{dominant process} has the largest summed duration and determines grouping; each worker's statistics still use all retained events. In Table~\ref{tab:metrics}, $\#$ counts elements, $\mathbf{1}$ is one when its condition holds and zero otherwise, and $\arg\max$ selects the process with the greatest summed duration.

\emph{Runs} group consecutive events of the same category within a recording. Their durations sum event durations, so gaps may be bridged and overlaps counted. We calculate each eligible worker's mean VA-run duration in minutes, then report the factory median. No eligible worker has an empty VA-run list; the implementation would return zero for that case.

\emph{VA transition rates} count changes from VA to another category in concatenated worker run sequences and divide by summed event-hours. This includes 15 transitions across clip joins among eligible workers; the counts do not establish interruptions in continuous work. As in the original script, worker VA shares, transition rates, and mean VA-run durations are rounded to one decimal before factory medians are calculated.

All per-video logs are nonempty, with no negative or zero durations. Overlaps count in duration totals and are merged only for coverage. Ambiguous transcript timestamps are used as recorded.

\subsection{Existing semantic audit protocol}
\label{sec:audit}
The audit compares extracted events with generated text using two samples: 42 events and 42 transcript entries, seven of each per factory. One author judged both samples. The record gives a random seed of 42 and sampling by factory, but not the full procedure, candidate pool, selection order, or treatment of timestamp formats. The sampling cannot be repeated exactly from these records.

For \emph{event support}, C means the overlapping transcript fully supports the label, P means the label is plausible but not fully supported, and X means contradiction. The sample contains 36 C events, six P events, and no X events. Strict support is $36/42=0.857$, with partial cases kept in the denominator.

For \emph{transcript-entry overlap}, COV means an entry overlaps at least one event, MISS means an uncovered entry judged operational, and EXC means an uncovered entry judged non-operational. There are 34 COV entries, four MISS entries, and four EXC entries. Keeping MISS and excluding EXC yields $34/(34+4)=0.895$, the \emph{reported entry-overlap rate}. It does not measure action recall: one COV entry describes camera tilt covered by a longer event, and the 38 retained entries are not independently verified as operational. Overlap alone does not establish complete or correct extraction of actions and labels.

The audit measures agreement with generated transcripts and cannot detect errors already present in them. There is no independent evaluator, inter-rater agreement estimate, or video audit. The sample rates do not bound errors in time shares, categories, observation selection, or causal interpretation. Appendix~\ref{supp-app:audit} of the Additional Material reproduces every judgment and note, including the flagged, unverified truncation explanation.
